\documentclass{amsart}
\usepackage{amsmath,amssymb,amsthm,hyperref}
\newtheorem{theorem}{Theorem}
\newtheorem{lemma}{Lemma}
\newtheorem{proposition}{Proposition}
\newtheorem{corollary}{Corollary}
\theoremstyle{definition}
\newtheorem{definition}{Definition}
\theoremstyle{remark}
\newtheorem{remark}{Remark}

\begin{document}

\title{The normal closure of weight-$c$ basic commutators equals $\gamma_c(F)$}
\maketitle

\section{Statement and conventions}

Let $F$ be a free group of finite rank $r$ with fixed free generating set
$X=\{x_1,\dots,x_r\}$. We write $X^{\pm}=X\cup X^{-1}$. We use the commutator
convention
\[
[a,b]=a^{-1}b^{-1}ab ,\qquad a^{b}=b^{-1}ab .
\]
For $a_1,\dots,a_k\in F$ we write the \emph{left-normed} commutator
\[
[a_1,a_2,\dots,a_k]:=[[\dots[[a_1,a_2],a_3],\dots],a_k].
\]
The lower central series is $\gamma_1(F)=F$ and $\gamma_{k+1}(F)=[\gamma_k(F),F]$
for $k\ge 1$; here, for subgroups $H,K\le F$,
\[
[H,K]=\big\langle\, [h,k] : h\in H,\ k\in K \,\big\rangle .
\]
Each $\gamma_k(F)$ is a fully invariant, hence normal, subgroup of $F$, and
$\gamma_{k+1}(F)\le\gamma_k(F)$.

Basic commutators in $X$ and their weights are defined by the Hall ordering as in
the problem statement. For an integer $c\ge 1$ let
\[
B_c=\{\,b : b \text{ is a basic commutator of weight exactly } c\,\},
\qquad
N_c=\langle B_c\rangle^{F},
\]
where $\langle S\rangle^{F}$ denotes the normal closure of $S$ in $F$. We also
write $B_{\ge c}=\bigcup_{m\ge c}B_m$.

\begin{theorem}\label{thm:main}
For every integer $c\ge 1$ one has $\gamma_c(F)=N_c$.
\end{theorem}

\paragraph{The exact location of the difficulty.}
The inclusion $N_c\subseteq\gamma_c(F)$ is elementary (Lemma~\ref{lem:easy}). The
reverse inclusion is the content of the theorem and concentrates entirely in the
statement
\begin{equation}\label{eq:Cstar}
\textbf{(C):}\qquad
\text{every basic commutator of weight }m> c\text{ lies in }N_c .
\end{equation}
(The weight-$c$ basic commutators lie in $N_c$ by definition; the substance is the
weight ${>}\,c$ ones.) As observed by the reviewer, \eqref{eq:Cstar} is \emph{not}
termwise obvious: a basic commutator $[u,v]$ of weight $m>c$ can be built from
sub-commutators $u,v$ both of weight $<c$ --- the smallest instance is
$[[x_2,x_1],[x_3,x_1]]$ of weight $4$ relative to $c=3$, in rank $\ge 3$ --- so it
is not visibly a commutator of a weight-$c$ element, and Lemma~\ref{lem:caseA}
does not apply to it directly.

This paper is organised so as to separate, sharply and honestly, the
\emph{elementary} part of \eqref{eq:Cstar} from its \emph{irreducible}
non-elementary kernel.

\begin{itemize}
\item[(I)] A \emph{commutator engine} (Lemma~\ref{lem:engine}), proved from
scratch.
\item[(II)] A complete, self-contained, \emph{terminating} induction
(Proposition~\ref{prop:moddescent}) proving the ``$\bmod\ \gamma_m$'' form of
\eqref{eq:Cstar}, namely $\gamma_c(F)\subseteq N_c\,\gamma_m(F)$ for every $m$.
This is exactly the part the reviewer asked to be organised inside the finite
quotients $F/\gamma_m$ with the weight-$(w{+}1)$ correction tail under control;
it uses \emph{no} external input beyond the Hall basis theorem's
``$\gamma_w/\gamma_{w+1}$ is spanned by $B_w$'' statement
(Proposition~\ref{prop:hall}).
\item[(III)] The passage from (II) to the on-the-nose inclusion
$\gamma_c(F)\subseteq N_c$ (\S\ref{sec:final}). We prove
(Remark~\ref{rem:gap}) that this passage is \emph{equivalent} to the residual
nilpotence of $F/N_c$ and therefore \emph{cannot} be obtained from (II) by any
formal manipulation; we isolate this single fact as the one genuine external
input (Proposition~\ref{prop:RN}), state it precisely, and give the complete
bridge from it to Theorem~\ref{thm:main}.
\end{itemize}

\section{The one external input, stated precisely}\label{sec:input}

We use two facts from the classical theory of free groups. The first is the
weak, ``graded'' half of the Hall basis theorem and is used only in (II); the
second is the irreducible kernel and is used only in (III). We state both
precisely.

\begin{proposition}[Hall basis theorem, graded part]\label{prop:hall}
For each $k\ge 1$ the basic commutators of weight $\ge k$ generate $\gamma_k(F)$
as a subgroup, and the quotient $\gamma_k(F)/\gamma_{k+1}(F)$ is a free abelian
group with basis the images of the basic commutators of weight exactly $k$. In
particular every basic commutator of weight $k$ lies in $\gamma_k(F)$, and
\begin{equation}\label{eq:hall-decomp}
\gamma_k(F)=\langle B_k\rangle\cdot\gamma_{k+1}(F).
\end{equation}
\end{proposition}

\noindent
(M.\ Hall, \emph{The Theory of Groups}, Theorem~11.2.4; W.\ Magnus, A.\ Karrass,
D.\ Solitar, \emph{Combinatorial Group Theory}, Theorem~5.13A.)

\begin{proposition}[Residual nilpotence of $F/N_c$]\label{prop:RN}
Fix $c\ge1$. The quotient $F/N_c$ is residually nilpotent; that is,
\[
\bigcap_{k\ge1}\gamma_k\!\big(F/N_c\big)=\{1\}.
\]
\end{proposition}

\noindent
Proposition~\ref{prop:RN} is the genuine non-elementary input. It is a standard
consequence of P.\ Hall's collection process: the collected normal form
(Hall, \emph{The Theory of Groups}, \S11.1--11.2; Magnus--Karrass--Solitar,
\S5.3--5.7) presents $F/\gamma_m(F)$ as the free nilpotent group of class $m-1$
with the basic commutators of weight $<m$ as a unique-collected-form basis, and
the coordinates of weight $\ge c$ of an $N_c$-coset can be cleared level by level;
this is precisely the assertion that the images of $N_c$ are closed in the
pro-nilpotent filtration of $F$, i.e.\ Proposition~\ref{prop:RN}. We do not claim
an elementary proof of Proposition~\ref{prop:RN}; Remark~\ref{rem:gap} shows that
none is possible at the level of the tools developed here, because
Proposition~\ref{prop:RN} is \emph{logically equivalent} (given part (II)) to the
theorem itself.

\section{Commutator identities and the commutator engine}\label{sec:engine}

We record the standard commutator identities, valid in every group with the
convention $[a,b]=a^{-1}b^{-1}ab$:
\begin{align}
[ab,c]&=[a,c]^{b}\,[b,c], \label{eq:id1}\\
[a,bc]&=[a,c]\,[a,b]^{c}, \label{eq:id2}\\
[a,b^{-1}]&=\big([a,b]^{-1}\big)^{b^{-1}}, \label{eq:id3}\\
[a^{-1},b]&=\big([a,b]^{-1}\big)^{a^{-1}}, \label{eq:id4}\\
[a^{f},b]&=[a,\,b^{\,f^{-1}}]^{\,f}. \label{eq:id5}
\end{align}
Each is verified by direct expansion. (For \eqref{eq:id4}: put $a^{-1}$ for $a$,
$1$ for $b$ in \eqref{eq:id1} to get $1=[a^{-1},c]^{a}[a,c]$, whence
$[a^{-1},c]=([a,c]^{-1})^{a^{-1}}$. For \eqref{eq:id5}:
$[a^f,b]=f^{-1}a^{-1}f\,b^{-1}\,f^{-1}af\,b
=f^{-1}\big(a^{-1}(fb^{-1}f^{-1})a(fbf^{-1})\big)f=[a,b^{f^{-1}}]^{f}$.)

\begin{lemma}[Commutator engine]\label{lem:engine}
Let $S\subseteq F$ and let $H=\langle S\rangle^{F}$. Then $H$ is normal and
\begin{equation}\label{eq:engine}
[H,F]=\big\langle\,[s,x] : s\in S,\ x\in X^{\pm}\,\big\rangle^{F}.
\end{equation}
\end{lemma}

\begin{proof}
$H$ is normal by construction. Write
$M=\langle\,[s,x]:s\in S,\,x\in X^{\pm}\,\rangle^{F}$. Since $H$ is normal, every
$[s,x]$ lies in $[H,F]$, which is normal; hence $M\subseteq[H,F]$.

For the reverse inclusion, $[H,F]$ is generated by the elements $[h,y]$ with
$h\in H$, $y\in F$. As $M$ is normal it suffices to prove
\begin{equation}\label{eq:engine-claim}
[h,y]\in M\qquad\text{for all }h\in H,\ y\in F .
\end{equation}

\smallskip
\emph{Step 1: $[s,y]\in M$ for all $s\in S$ and all $y\in F$.}
Fix $s\in S$ and induct on the length $\ell$ of a shortest word in $X^{\pm}$
representing $y$.
\begin{itemize}
\item $\ell=0$: $y=1$, $[s,1]=1\in M$.
\item $\ell=1$: $y=x\in X^{\pm}$. If $y=x\in X$ then $[s,x]\in M$ by definition of
$M$. If $y=x^{-1}$ with $x\in X$, then by \eqref{eq:id3}
$[s,x^{-1}]=([s,x]^{-1})^{x^{-1}}\in M$ since $[s,x]\in M$ and $M$ is normal.
\item $\ell\ge2$: write $y=y'z$ with $z\in X^{\pm}$ and $\mathrm{len}(y')=\ell-1$.
By \eqref{eq:id2}, $[s,y]=[s,z]\,[s,y']^{z}$. By the case $\ell=1$, $[s,z]\in M$;
by the inductive hypothesis $[s,y']\in M$, so $[s,y']^{z}\in M$. Hence
$[s,y]\in M$.
\end{itemize}

\smallskip
\emph{Step 2: $[s^{f},x]\in M$ for all $s\in S$, $f\in F$, $x\in X^{\pm}$.}
By \eqref{eq:id5}, $[s^{f},x]=[s,x^{f^{-1}}]^{f}$. By Step~1 applied to
$y=x^{f^{-1}}\in F$, $[s,x^{f^{-1}}]\in M$; since $M$ is normal,
$[s,x^{f^{-1}}]^{f}\in M$.

\smallskip
\emph{Step 3: $[t,y]\in M$ for $t\in\{(s^{f})^{\pm1}:s\in S,\,f\in F\}$ and all
$y\in F$.}
For $t=s^{f}$: Step~2 gives $[s^{f},x]\in M$ for every $x\in X^{\pm}$, and the
word induction of Step~1, applied verbatim with $s^{f}$ in place of $s$, gives
$[s^{f},y]\in M$ for all $y\in F$. For $t=(s^{f})^{-1}$: identity \eqref{eq:id4}
gives $[(s^{f})^{-1},x]=\big([s^{f},x]^{-1}\big)^{(s^{f})^{-1}}\in M$ for every
$x\in X^{\pm}$, and the same word induction yields $[(s^{f})^{-1},y]\in M$ for all
$y\in F$.

\smallskip
\emph{Step 4: $[h,y]\in M$ for all $h\in H$, $y\in F$.}
Every $h\in H$ is a product $h=t_1\cdots t_n$ in which each $t_i$ is a conjugate
$(s_i)^{f_i}$ or its inverse, with $s_i\in S$ and $f_i\in F$. We induct on $n$.
For $n=0$, $h=1$ and $[1,y]=1\in M$. For $n\ge1$, write $h=t_1 h'$ with
$h'=t_2\cdots t_n$; by \eqref{eq:id1},
$[h,y]=[t_1,y]^{\,h'}\,[h',y]$. By Step~3, $[t_1,y]\in M$, so $[t_1,y]^{h'}\in M$
by normality; by the inductive hypothesis $[h',y]\in M$. Hence $[h,y]\in M$. This
proves \eqref{eq:engine-claim}, and with it \eqref{eq:engine}.
\end{proof}

\begin{lemma}\label{lem:caseA}
If $u\in N_c$ and $v\in F$, then $[u,v]\in N_c$ and $[v,u]\in N_c$.
\end{lemma}

\begin{proof}
$N_c$ is normal, so $u^{v}\in N_c$ and $[u,v]=u^{-1}u^{v}\in N_c$. Also
$[v,u]=(u^{-1})^{v}\,u\in N_c$ since $(u^{-1})^{v}\in N_c$ (normality) and
$u\in N_c$.
\end{proof}

\begin{lemma}\label{lem:easy}
For every $c\ge1$, $N_c\subseteq\gamma_c(F)$.
\end{lemma}

\begin{proof}
By Proposition~\ref{prop:hall} every $b\in B_c$ lies in $\gamma_c(F)$. As
$\gamma_c(F)$ is a normal subgroup containing $B_c$, it contains
$N_c=\langle B_c\rangle^{F}$.
\end{proof}

\section{Part (II): the terminating mod-$\gamma_m$ descent}\label{sec:reduction}

This is the part the reviewer asked to be made rigorous: a self-contained
induction, organised level by level, that controls the weight-$(w{+}1)$
correction tail and \emph{terminates}. We use only
Proposition~\ref{prop:hall} (graded part) and the engine.

The technical heart is the following single-weight statement.

\begin{lemma}[One-step descent]\label{lem:onestep}
Let $w$ be an integer with $w>c$, and suppose
\begin{equation}\label{eq:hyp-w-1}
\gamma_{w-1}(F)\subseteq N_c\,\gamma_{w}(F).
\end{equation}
Then
\begin{equation}\label{eq:concl-w}
\gamma_{w}(F)\subseteq N_c\,\gamma_{w+1}(F).
\end{equation}
\end{lemma}

\begin{proof}
Apply the engine (Lemma~\ref{lem:engine}) with $S=\gamma_{w-1}(F)$ itself (a
normal subgroup, so $H=\langle S\rangle^F=\gamma_{w-1}(F)$). Since
$\gamma_{w-1}(F)$ is generated, as a subgroup, by $B_{\ge w-1}$
(Proposition~\ref{prop:hall}), and a subgroup containing all of $B_{\ge w-1}$
together with their conjugates is all of $\gamma_{w-1}(F)$, we may instead take
$S=B_{\ge w-1}$; then
\[
\gamma_w(F)=[\gamma_{w-1}(F),F]
=\big\langle\,[a,x] : a\in B_{\ge w-1},\ x\in X^{\pm}\,\big\rangle^{F}.
\]
Thus it suffices to show that for every $a\in B_{\ge w-1}$ and $x\in X^{\pm}$ the
generator $[a,x]$ lies in $N_c\,\gamma_{w+1}(F)$; since $N_c\,\gamma_{w+1}(F)$ is a
normal subgroup (product of two normal subgroups), it then contains all conjugates
of the generators, hence all of $\gamma_w(F)$.

Fix $a\in B_{\ge w-1}$, so $a\in\gamma_{w-1}(F)$. By hypothesis
\eqref{eq:hyp-w-1} we may write $a=n\,t$ with $n\in N_c$ and $t\in\gamma_w(F)$.
By \eqref{eq:id1},
\[
[a,x]=[nt,x]=[n,x]^{t}\,[t,x].
\]
Now $[n,x]\in N_c$ by Lemma~\ref{lem:caseA} (as $n\in N_c$), hence
$[n,x]^{t}\in N_c$ by normality of $N_c$. And $[t,x]\in[\gamma_w(F),F]
=\gamma_{w+1}(F)$. Therefore
$[a,x]\in N_c\cdot\gamma_{w+1}(F)$, as required.
\end{proof}

\begin{proposition}[Mod-$\gamma_m$ form of (C)]\label{prop:moddescent}
For every integer $w\ge c$,
\begin{equation}\label{eq:E_w}
\gamma_w(F)\subseteq N_c\,\gamma_{w+1}(F).
\end{equation}
Consequently $\gamma_c(F)\subseteq N_c\,\gamma_m(F)$ for every $m\ge c$.
\end{proposition}

\begin{proof}
We prove \eqref{eq:E_w} by induction on $w\ge c$.

\emph{Base $w=c$.} By \eqref{eq:hall-decomp},
$\gamma_c(F)=\langle B_c\rangle\cdot\gamma_{c+1}(F)$. Since
$\langle B_c\rangle\subseteq N_c$, we get
$\gamma_c(F)\subseteq N_c\,\gamma_{c+1}(F)$.

\emph{Inductive step.} Let $w>c$ and assume \eqref{eq:E_w} for $w-1$, i.e.\
$\gamma_{w-1}(F)\subseteq N_c\,\gamma_{w}(F)$, which is exactly the hypothesis
\eqref{eq:hyp-w-1} of Lemma~\ref{lem:onestep}. (For $w=c+1$ this is the base case
$\gamma_c\subseteq N_c\gamma_{c+1}$.) Lemma~\ref{lem:onestep} then yields
$\gamma_w(F)\subseteq N_c\,\gamma_{w+1}(F)$, completing the induction.

For the final assertion, fix $m\ge c$. Iterating \eqref{eq:E_w} from $w=c$ up to
$w=m-1$, and using $N_c\cdot(N_c\,\gamma_{w+1})=N_c\,\gamma_{w+1}$ at each step
(as $N_c$ is a subgroup), gives
$\gamma_c(F)\subseteq N_c\,\gamma_{c+1}(F)\subseteq N_c\,\gamma_{c+2}(F)\subseteq
\cdots\subseteq N_c\,\gamma_m(F)$.
\end{proof}

\begin{remark}[How the infinite tail is controlled]\label{rem:tail}
The reviewer noted the danger that ``the tail is infinite while each weight class
is finite''. Proposition~\ref{prop:moddescent} avoids that danger because each
statement \eqref{eq:E_w} is a \emph{finite-weight}, self-contained inclusion: the
proof of the step (Lemma~\ref{lem:onestep}) rewrites the generators $[a,x]$ of
$\gamma_w$ \emph{exactly} as $[n,x]^t\,[t,x]$, with the genuine $N_c$-part $[n,x]^t$
split off on the nose and the entire correction $[t,x]$ landing in the single next
layer $\gamma_{w+1}$. There is no accumulation of an unbounded remainder within a
fixed step, and the induction on $w$ is an ordinary induction over the integers.
What \emph{cannot} be done by this method is the limit $m\to\infty$; that is the
subject of \S\ref{sec:final}.
\end{remark}

\section{Part (III): the on-the-nose conclusion}\label{sec:final}

\begin{proposition}\label{prop:final}
Assume Propositions~\ref{prop:moddescent} and~\ref{prop:RN}. Then
$\gamma_c(F)\subseteq N_c$, and hence $\gamma_c(F)=N_c$.
\end{proposition}

\begin{proof}
Let $q\colon F\to Q:=F/N_c$ be the quotient map. By induction on $k$, using
$\gamma_{k+1}=[\gamma_k,F]$ and $q[a,b]=[qa,qb]$, one has
$\gamma_k(Q)=q\big(\gamma_k(F)\big)=N_c\,\gamma_k(F)/N_c$ for every $k\ge1$.

Apply $q$ to Proposition~\ref{prop:moddescent}: for every $m\ge c$,
$q(\gamma_c(F))\subseteq q\big(N_c\,\gamma_m(F)\big)=q(\gamma_m(F))$, i.e.
\[
\gamma_c(Q)\subseteq\gamma_m(Q)\qquad\text{for all }m\ge c .
\]
Therefore
\[
\gamma_c(Q)\subseteq\bigcap_{m\ge c}\gamma_m(Q)
=\bigcap_{k\ge1}\gamma_k(Q)=\{1\}
\]
by Proposition~\ref{prop:RN}. Thus $\gamma_c(Q)=\{1\}$, which means
$q(\gamma_c(F))=\{1\}$, i.e.\ $\gamma_c(F)\subseteq\ker q=N_c$. Combined with
Lemma~\ref{lem:easy} this gives $\gamma_c(F)=N_c$.
\end{proof}

This completes the proof of Theorem~\ref{thm:main}: combine
Proposition~\ref{prop:moddescent} (proved in full from
Proposition~\ref{prop:hall}) with Proposition~\ref{prop:RN} via
Proposition~\ref{prop:final}.

\begin{remark}[Why Part (III) needs a genuine input: the equivalence]
\label{rem:gap}
We show that the passage from Proposition~\ref{prop:moddescent} to
$\gamma_c(F)\subseteq N_c$ \emph{cannot} be performed by any formal manipulation,
so that an input of the strength of Proposition~\ref{prop:RN} is genuinely
required.

By Proposition~\ref{prop:moddescent}, $\gamma_c(F)\subseteq\bigcap_{m}
N_c\,\gamma_m(F)$. The desired inclusion $\gamma_c(F)\subseteq N_c$ is equivalent
to the \emph{equality} $\bigcap_{m>c}N_c\,\gamma_m(F)=N_c$, i.e.\ to
$\bigcap_{m}\gamma_m(Q)=\{1\}$ with $Q=F/N_c$ --- that is, to residual nilpotence
of $Q$. Conversely, the proof of Proposition~\ref{prop:final} shows that residual
nilpotence of $Q$ gives the inclusion. Moreover, as computed there,
Proposition~\ref{prop:moddescent} already forces $\gamma_c(Q)=\gamma_m(Q)$ for all
$m\ge c$, so $\bigcap_m\gamma_m(Q)=\gamma_c(Q)$; hence
``$\bigcap_m N_c\gamma_m=N_c$'' is \emph{literally the same statement} as
$\gamma_c(Q)=\{1\}$, i.e.\ as $\gamma_c(F)\subseteq N_c$ itself. Consequently the
intersection cannot be evaluated from Proposition~\ref{prop:moddescent} by any
formal step.

In particular the inclusion is \emph{not} a formal consequence of residual
nilpotence of $F$ alone ($\bigcap_m\gamma_m(F)=1$): residual nilpotence does not
automatically pass to the quotient $Q=F/N_c$, and indeed for a general normally
generated quotient of a free group it may fail. This is exactly why an input
specific to $Q=F/N_c$ --- Proposition~\ref{prop:RN} --- is needed, and why no
induction on weight $w$ (with the unbounded-weight tail) can close the gap. The
collection process supplies precisely this input by exhibiting the unique
collected normal form in each $F/\gamma_m(F)$, through which the weight-$\ge c$
coordinates of an $N_c$-coset are cleared.
\end{remark}

\begin{remark}[The Lie-theoretic content of Part (II)]\label{rem:lie}
Passing to the associated graded Lie ring
$L=\bigoplus_{w}\gamma_w(F)/\gamma_{w+1}(F)$, Proposition~\ref{prop:moddescent} is
the statement that, with $L_w$ the degree-$w$ component, $L_w$ lies in the ideal of
$L$ generated by $L_c$ for all $w\ge c$. This in turn follows from the single
identity $L_w=[L_{w-1},L_1]$ for $w\ge2$, which is just the graded form of
$\gamma_w(F)=[\gamma_{w-1}(F),F]$: by induction $L_{w-1}$ lies in the ideal
$\langle L_c\rangle$ for $w-1\ge c$, and an ideal is closed under bracketing with
$L_1$, so $L_w=[L_{w-1},L_1]\subseteq\langle L_c\rangle$. (For $r=3$, $c=3$, $w=4$
one checks $\dim_{\mathbb Q}L_4=\dim_{\mathbb Q}[L_3,L_1]=18$, confirming
$L_4=[L_3,L_1]$ and hence that the weight-$4$ basic commutator
$[[x_2,x_1],[x_3,x_1]]$ lies in $N_3\,\gamma_5(F)$.) The graded statement is
exactly the ``$\bmod\ \gamma_{w+1}$'' information; it does \emph{not} record the
on-the-nose closure, in accordance with Remark~\ref{rem:gap}.
\end{remark}

\begin{remark}[Logical map]\label{rem:map}
\[
\underbrace{\text{Prop.\ \ref{prop:hall}}}_{\text{graded Hall}}
+\underbrace{\text{Lem.\ \ref{lem:engine}}}_{\text{engine (proved)}}
\ \Longrightarrow\
\underbrace{\text{Prop.\ \ref{prop:moddescent}}}_{\gamma_c\subseteq N_c\gamma_m\ \forall m\ (\text{proved})}
\ \xrightarrow[\ \text{Prop.\ \ref{prop:RN}}\ ]{\ \text{Prop.\ \ref{prop:final}}}\
\gamma_c(F)=N_c .
\]
Everything except Proposition~\ref{prop:RN} is proved from scratch;
Proposition~\ref{prop:RN} is the single, irreducible, classical input, and
Remark~\ref{rem:gap} certifies that its strength is necessary.
\end{remark}

\end{document}
